\scanpage{57}
$H^1(S,T_1)$ \unsure{classification des formes intérieures munies} d'un couple de Killing.

[N.B. si $G$ est une forme intérieure, \unsure{alors} deux couples de Killing \inserted{$(T,B),(T',B')$} dans $G$ définissent le $=$ élément de $H^1(S,T_1)$ ssi ils sont conjugués par un \unsure{aut.} intérieur. Car \unsure{à priori leur} \unsure{différence pour une certaine classe} dans $H^1(S,\unsure{\overline B|\overline T})$ est \unsure{valable} --- si $T=T'$, c'est la classe \inserted{mod $\unsure D$} de l'unique $\sigma\in W(T)$ qui transforme $B$ en $B'$; si $B'$ est l'opposé de $B$ rel. à $T$, alors le $\sigma$ en question est \deleted{\unsure{donné par}} \unsure{l'image d'un} élément \unsure{connu} de $H^1(S,T)$]

\margin{$H^1(S,T)\mathbin{-}H^1(S,B)$}

La question de savoir si deux couples de Killing dans un $=$ $G$ sont tjrs conjugués, \unsure{équivaut} donc à celle de savoir si $H^1(S,T_1)\mathbin{-}H^1(S,G_1)$ est tjrs injective. Pour un contre-exemple, il suffit de montrer que $H^1(S,T_1)\mathbin{-}H^1(S,N_1)$ n'est pas tjrs injective. Or deux \unsure{élémts} de $H^1(S,T_1)$ ont la $=$ image dans $H^1(S,N_1)$ ssi ils sont conjugués par $H^0(S,W_1)$ opérant de façon évidente. Je prendrai \unsure{$G$ pour} $G_0=\unsure{GL|GP}(2)$, avec $E=\mathbf Z/2\mathbf Z$, $W_0=\mathbf Z/2\mathbf Z$ est commutatif, \deleted{\unsure{donc} $\unsure{W_1=W_0}$ \unsure{comme}} $E$ y opère trivialt (il n'a pas le choix!) donc $W_0=W_1$, qui opère par symétrie sur $T_1$.
\margin{56}
